\section{Finite convex geometry and Euler's identity}\label{sec:euler}

We first record the finite geometric facts used in the proof. A line
intersects a convex polygon in an interval: substitute its affine
parameter into the finitely many half-plane inequalities and take the
largest lower bound and the smallest upper bound. A convex polygon is
triangulated by the diagonals from any one vertex. A point in the polygon
is located in one of those triangles by finitely many determinant signs,
and its barycentric coefficients are obtained by a $2\times2$ inverse.
These statements use exact arithmetic and finite minimum selection.

At a vertex $v$ of a three-dimensional convex polyhedron, intersect its
tangent cone with a plane strictly separating $v$ from the other
vertices. Such a plane is explicit: add the inward deficits of the
incident supporting planes, which vanish simultaneously only at $v$,
and choose a positive level smaller than all the finitely many values
at other vertices. The section is a convex polygon. Its vertices are
the intersections with the edge rays, in their cyclic order. A direction
from $v$ into the polyhedron is a nonnegative combination of edge rays:
intersect it with the section and use the preceding polygon triangulation.
Inserting diagonals in facets merely subdivides edges of this vertex
figure. These observations prove the local cyclic-link properties used
below directly from the finite supporting inequalities.

\begin{lemma}[Generic finite height]\label{lem:height}
For $V$ distinct exactly represented points in three dimensions one can
construct a linear functional taking distinct values on all the points
by trying at most $2\binom V2+1$ integer candidates.
\end{lemma}
\begin{proof}
Try $h_t(x,y,z)=x+ty+t^2z$ for integers
$0\leq t\leq2\binom V2$. For any distinct pair of points, equality of
their heights is a nonzero polynomial equation of degree at most two
in $t$. Such a polynomial has at most two distinct roots, by division
by linear factors. There are at most $2\binom V2$ forbidden values.
Testing the displayed list therefore returns a valid candidate.
This is the finite pigeonhole argument on decidable lists.
\end{proof}

\begin{lemma}[Euler by lower-neighbor counting]\label{lem:euler}
For a finite full-dimensional convex polyhedron,
\[
 V+F=E+2.
\]
\end{lemma}
\begin{proof}
Triangulate each facet by a fan without adding vertices. Each diagonal
adds one edge and one face, so the value of $V-E+F$ is unchanged.
Choose the generic height of Lemma~\ref{lem:height}. Assign every edge
and every triangle to its highest vertex. At a vertex $v$ let $d_v$
count its lower neighbors and $t_v$ its incident triangles whose other
two vertices are lower. Its contribution to $V-E+F$ is
\[
 1-d_v+t_v.
\]
In the convex vertex figure, being lower is an affine half-plane
condition. The lower neighbors therefore form a consecutive block
in the cyclic order, unless all or none of them are lower. For a
nonempty proper block, $t_v=d_v-1$, and the contribution is zero.
If all neighbors are lower, the edge-ray combination property proves
that $v$ is the unique global maximum; its full cyclic link gives
$t_v=d_v$, so its contribution is one. If none is lower, the same
argument makes $v$ the unique global minimum; its contribution is one.
There are no other local extrema. Summing gives two. Remove the added
facet diagonals to obtain the stated identity.
\end{proof}

No spherical topology theorem has been used here. For an internal
arithmetic proof write the contribution calculation as an equality of
two accumulated natural-number counts, avoiding truncated subtraction.

\subsection{A finite planar drawing}

Some later arguments use triangular regions in a planar drawing.
Here is a construction of that drawing, including the separation it
requires. Choose a facet $F_0$ and average its vertices to obtain $c$.
This average lies strictly inside every other facet half-space: another
facet cannot contain all of $F_0$, and every coefficient in the average
is positive. Move a small distance outward from $c$ across $F_0$,
obtaining a point $o$ still strictly inside all the other half-spaces.
The distance can be chosen by taking the minimum of their positive
slacks divided by twice a bound on the change of their linear forms.

We may additionally require that $o$ lie in no plane through any
noncollinear triple of vertices. To see that this is a bounded choice,
take an exact box around the first choice still satisfying all its
strict inequalities. If there are $m$ forbidden planes, a Cartesian
grid of $m+1$ values in each coordinate has $(m+1)^3$ points. Each
plane meets at most $(m+1)^2$ grid points, so at least one avoids all
the planes. All tests are finite exact determinant comparisons.

Project centrally from $o$ onto the plane of $F_0$. Every vertex not on
$F_0$ projects strictly inside $F_0$: its projection is a proper convex combination
of that vertex and $o$, and hence satisfies all the other supporting
inequalities strictly. For a ray through a point of $F_0$, intersection
with the polyhedron is an interval. Its far endpoint is the minimum
of the finitely many exit parameters supplied by the facet inequalities.
Thus the remaining facets tile $F_0$ under projection, with disjoint
relative interiors and no crossing edges. This is a Schlegel drawing,
obtained without compactness, an intermediate value theorem, or a
topological sphere axiom.

In a triangular-face polyhedron every graph three-cycle is a
nondegenerate geometric triangle in this generic drawing. Its inside
is given by its three half-plane tests. No edge can cross from that
inside to the outside without meeting the boundary. Consequently a
three-cycle which is not a facet has vertices on both sides. If a
side had no vertex, the triangle on that side could have neither an
internal edge nor a subdivision and would be a facet. Either side
can be made the bounded triangle by choosing a facet on the other
side as the outer facet. This requires only the separation property
of an ordinary triangle, not a general Jordan curve theorem.

\begin{lemma}[Triangular disk count]\label{lem:disk}
Suppose a geometric triangle is tiled by a finite straight-line
triangulation with $k$ interior vertices and no extra boundary vertices.
Let $e_{II}$ and $e_{IB}$ count edges with respectively two and one
interior endpoints. If every interior vertex has degree five in the
triangulation, then $e_{II}=2k$.
\end{lemma}
\begin{proof}
Write $E_D,T_D$ for all edges and small triangles. A generic planar
height gives
\[
 (k+3)-E_D+T_D=1.
\]
Indeed, use the same assignment to highest vertices as before.
An interior vertex has a nonempty proper block of lower neighbors
in its circular link, contributing zero. At a boundary vertex the
link is a path in the convex triangular corner. All vertices except
the unique minimum have a nonempty consecutive block of lower
neighbors and contribute $1-d+(d-1)=0$; the minimum contributes one.
Counting triangle-edge incidences also gives $3T_D=2E_D-3$.
Eliminating $T_D$ yields $E_D=3k+3$, so
\[
 e_{II}+e_{IB}=3k.
\]
Counting edges incident to the $k$ interior vertices gives
$2e_{II}+e_{IB}=5k$. Subtract the preceding equality, or cancel common
summands in natural numbers, to obtain $e_{II}=2k$.
\end{proof}

\section{The numerical classification}\label{sec:numerical}

Every edge has two endpoints and belongs to exactly two faces. Counting
vertex-edge and face-edge incidences therefore gives
\begin{equation}\label{eq:incidence}
 qV=2E=pF.
\end{equation}
There is at least one edge, and $p,q\geq3$. Put
$p=a+2$, $q=b+2$, so $a,b\geq1$. Multiply Euler's identity by $pq$
and use~\eqref{eq:incidence}. Expanding and cancelling gives the
subtraction-free identity
\begin{equation}\label{eq:core}
 Eab+2pq=4E.
\end{equation}
If $ab\geq4$, then $Eab\geq4E$, contradicting the positive summand
$2pq$. Hence
\begin{equation}\label{eq:small}
 (p-2)(q-2)=ab<4.
\end{equation}
For positive integers $a,b$, either $a=1$ and $b=1,2,3$, or
$b=1$ and $a=2,3$. If both were at least two their product would be
at least four. These are decidable arithmetic cases, not a
nonconstructive dichotomy.

Substitution into~\eqref{eq:core} and~\eqref{eq:incidence} gives the
complete table:
\begin{center}
\begin{tabular}{lrrrrr}
\toprule
Name & $p$ & $q$ & $V$ & $E$ & $F$\\
\midrule
Tetrahedron & 3 & 3 & 4 & 6 & 4\\
Cube & 4 & 3 & 8 & 12 & 6\\
Octahedron & 3 & 4 & 6 & 12 & 8\\
Dodecahedron & 5 & 3 & 20 & 30 & 12\\
Icosahedron & 3 & 5 & 12 & 30 & 20\\
\bottomrule
\end{tabular}
\end{center}
At this point the names are labels for possible cases. Existence and
uniqueness have not been inferred just from this table. In particular,
the calculation alone would not exclude two different solids with the
same $(p,q)$. The next two sections supply both missing parts.
